\section{Arquitectura del sistema}
\label{sec:architecture}

\subsection{Visi\'on general de la pipeline}

La Figura~\ref{fig:pipeline} muestra las cinco etapas de la
pipeline. Una muestra se analiza est\'atica y din\'amicamente en
paralelo; ambas salidas alimentan un \emph{normalizador} que
deduplica IOCs y agrupa las observaciones ATT\&CK por t\'ecnica; un
\emph{mapeador} reconcilia la confianza entre fuentes y produce la
lista final de t\'ecnicas; un \emph{exportador} convierte esa lista en
un paquete STIX~2.1 y/o un evento MISP en vivo. Cada etapa es una
herramienta de l\'inea de comandos independiente que opera sobre un
esquema JSON bien definido, de modo que la pipeline puede ejecutarse
de principio a fin o volver a ejecutar cualquier etapa de forma
aislada, propiedad utilizada extensamente durante la validaci\'on: un
solo cambio en una regla de detecci\'on a menudo solo requiere volver a
ejecutar el mapeador, no la detonaci\'on completa en sandbox.

\begin{center}
\begin{tabular}{c}
\texttt{muestra} \\
$\downarrow$ \\
\texttt{An\'alisis est\'atico} \quad + \quad \texttt{An\'alisis din\'amico (CAPE)} \\
$\downarrow$ \\
\texttt{Normalizador} $\rightarrow$ \texttt{normalized\_iocs.json} \\
$\downarrow$ \\
\texttt{Mapeador (reconciliaci\'on de confianza)} $\rightarrow$ \texttt{attck\_mapping.json} \\
$\downarrow$ \\
\texttt{Exportador} $\rightarrow$ \texttt{paquete STIX 2.1} / \texttt{evento MISP}
\end{tabular}
\captionof{figure}{Etapas de la pipeline y artefactos intermedios.}
\label{fig:pipeline}
\end{center}

\subsection{Motor de an\'alisis est\'atico}
\label{sec:static-rules}

El motor est\'atico analiza binarios PE (nativos y .NET) y ELF,
extrayendo importaciones, cadenas (est\'andar y decodificadas por
FLOSS~\cite{floss}), secciones, coincidencias YARA y caracter\'isticas
de entrop\'ia. Para muestras empaquetadas como instalador (NSIS, Inno,
MSI) realiza una extracci\'on en varias pasadas: se extrae cada
fichero incluido, las claves se descubren en \emph{todos} los ficheros
extra\'idos antes de intentar ning\'un descifrado (de modo que una
clave incrustada en un componente pueda descifrar una carga en otro),
las claves candidatas se prueban contra cualquier fichero que parezca
cifrado, y todos los ficheros, descifrados o no, se analizan
completamente despu\'es. La \S\ref{sec:case-installer} detalla un error
real encontrado y corregido en c\'omo funcionaba antes esta
agregaci\'on.

Las reglas de detecci\'on mapean evidencia est\'atica a t\'ecnicas
ATT\&CK mediante tres mecanismos complementarios:

\begin{itemize}
  \item \textbf{Tablas de importaciones/cadenas.} Un mapeo curado de
    nombres de API de Windows o patrones de cadena espec\'ificos a
    identificadores de t\'ecnica.
  \item \textbf{Comprobaciones de combinaci\'on.} Para t\'ecnicas cuyo
    indicador individual m\'as fuerte tambi\'en es com\'un en software
    leg\'itimo, una regla exige una \emph{combinaci\'on coherente}
    espec\'ifica de indicadores antes de reportar confianza elevada. La
    Tabla~\ref{tab:combos} enumera todas las comprobaciones de
    combinaci\'on implementadas mediante coincidencia expl\'icita de
    m\'ultiples importaciones; el mismo mecanismo de recuento de
    corroboraci\'on tambi\'en rige varias reglas de patrones de cadena
    a\~nadidas durante la ampliaci\'on de cobertura (p.\ ej.\ el
    \code{mshta.exe} + \code{.hta} de Mshta, o los tres nombres de clase
    corroborantes de WMI Event Subscription), tabuladas en su lugar en
    el Ap\'endice~\ref{app:rules}.
  \item \textbf{Llamadas a la BCL de .NET.} La tabla de importaciones
    nativas de un binario .NET suele ser \'unicamente el stub de
    arranque del CLR, as\'i que los cuerpos de m\'etodos gestionados se
    escanean directamente en busca de llamadas a APIs espec\'ificas de
    la Base Class Library (p.\ ej.\ \code{Convert::FromBase64String}).
\end{itemize}

\begin{table}[h]
\centering
\scriptsize
\begin{tabular}{@{}p{3.0cm}p{6.6cm}p{3.0cm}@{}}
\toprule
\textbf{T\'ecnica} & \textbf{Combinaci\'on requerida} & \textbf{Rechazado en solitario} \\
\midrule
T1055 (Process Injection) & \code{VirtualAllocEx} + \code{WriteProcessMemory} + \code{CreateRemoteThread}, o el tr\'io de lectura \code{OpenProcess} + \code{VirtualQueryEx} + \code{ReadProcessMemory} & \code{OpenProcess} en solitario \\
\rowline
T1055.012 (Process Hollowing) & \code{CreateProcessW} + \code{WriteProcessMemory} + \code{SetThreadContext} + \code{ResumeThread} + \code{Nt/ZwUnmapViewOfSection} & cualquier API individual \\
\rowline
T1685 (Disable or Modify Tools) & \code{WTSEnumerateProcessesW} + \code{OpenProcess} + \code{TerminateProcess} & \code{TerminateProcess} sola \\
\rowline
T1010 (Application Window Discovery) & \code{EnumWindows} + \code{GetWindowTextW} & \code{GetWindowTextW} sola \\
\rowline
T1134.001/.003 (Impersonation/Theft; Make and Impersonate Token) & \code{ImpersonateLoggedOnUser} + \code{DuplicateTokenEx} (.001) o \code{LogonUserW} (.003) & \code{Impersonate\-LoggedOnUser} sola \\
\bottomrule
\end{tabular}
\caption{Reglas de detecci\'on sensibles a combinaci\'on. Un
componente listado como ``rechazado en solitario'' se excluye
deliberadamente de la promoci\'on de confianza por importaci\'on \'unica
porque tiene un uso leg\'itimo sustancial fuera de la t\'ecnica
afirmada.}
\label{tab:combos}
\end{table}

\subsection{Motor de an\'alisis din\'amico}
\label{sec:dynamic}

La detonaci\'on din\'amica se ejecuta dentro de CAPE~\cite{cape} contra
una red simulada con INetSim, aislada y sin salida real a Internet, un
dise\~no de contenci\'on prioritaria cuya consecuencia para la
observabilidad se discute en \S\ref{sec:limitations}. El informe en
bruto de CAPE se depura en un JSON m\'as peque\~no, centrado en IOCs:
firmas, ficheros soltados, actividad de red y una lista curada de
mapeos ATT\&CK.

Las firmas de la comunidad de CAPE llevan sus propias etiquetas de
t\'ecnica, pero se trataron como una hip\'otesis, no como datos de
referencia, durante toda la validaci\'on. La capa de curaci\'on
mantiene tres tablas de correcci\'on, cada entrada justificada
individualmente frente a la evidencia en bruto de la propia firma (su
campo \code{data}, no solo su nombre o categor\'ia):

\begin{itemize}
  \item \textbf{Firmas poco fiables}, descartadas por completo: p.\
    ej.\ \code{unbacked\_process\_mitigation\_alteration} (etiqueta
    \technique{T1562}), cuyas direcciones de llamada en bruto en la
    muestra que la activ\'o se sit\'uan todas en un rango estrecho de
    memoria alta consistente con una DLL del sistema, lo que sugiere
    que la clasificaci\'on de ``memoria sin respaldo'' fall\'o sobre la
    propia inicializaci\'on autom\'atica de procesos de Windows en lugar
    de sobre la acci\'on de la propia muestra.
  \item \textbf{Eliminaciones de t\'ecnica por firma}, donde una
    etiqueta en una firma por lo dem\'as fiable es incorrecta: p.\
    ej.\ \code{antiav\_servicestop} etiqueta \technique{T1489} (Service
    Stop), pero en la muestra que la activ\'o, sus propios datos en
    bruto nombran el servicio concreto detenido: el propio driver de la
    muestra, no el de un producto de seguridad.
  \item \textbf{Reasignaciones de t\'ecnica}, donde las propias
    definiciones de MITRE contradicen la etiqueta de la comunidad: p.\
    ej.\ \code{unbacked\_api\_resolution} est\'a etiquetada por la
    comunidad como \technique{T1129} (Shared Modules), que la propia
    p\'agina de MITRE define como un m\'odulo respaldado por un fichero
    en disco; la descripci\'on de la propia firma (``memoria asignada
    din\'amicamente, sin respaldo'') es exactamente lo contrario, y
    \technique{T1620} (Reflective Code Loading) encaja en su lugar. La
    misma tabla tambi\'en salva una diferencia de versi\'on: las firmas
    de la comunidad de CAPE quedan fijadas a la versi\'on de ATT\&CK
    con la que se escribieron por \'ultima vez, as\'i que las etiquetas
    \technique{T1562}/\technique{T1562.001} y \technique{T1070.001}
    que dispara CAPE se reasignan a sus identidades de ATT\&CK v19
    (\technique{T1685}, \technique{T1685.005}) en el momento de la
    curaci\'on, con independencia de la versi\'on propia de este
    proyecto (\S\ref{sec:limitations}).
\end{itemize}

Una capacidad adicional, a\~nadida durante la ampliaci\'on de
cobertura, comprueba los \code{executed\_commands} en bruto de CAPE
frente a patrones de comandos LOLBin/reconocimiento conocidos que no
tienen ninguna firma de comunidad correspondiente: p.\ ej.\
\code{netsh wlan show profiles}, que CAPE observa y registra pero no
mapea por s\'i mismo a ninguna t\'ecnica.

\subsection{Reconciliaci\'on de confianza}
\label{sec:reconciliation}

La regla central de reconciliaci\'on, implementada en
\code{reconcile.py}, es: una t\'ecnica observada tanto
por evidencia est\'atica como din\'amica alcanza confianza
\emph{alta}, salvo que la mejor confianza propia de ambas fuentes sea
\emph{baja}, en cuyo caso queda en \emph{media}; una t\'ecnica
observada por una sola fuente mantiene el nivel propio de esa fuente,
y no puede promocionarse por volumen de evidencia de una \'unica
fuente.

\subsubsection{Extensi\'on entre niveles}

Esta regla, tal como se implement\'o originalmente, comparaba
observaciones agrupadas por identificador de t\'ecnica \emph{exacto}.
En la pr\'actica, una firma din\'amica frecuentemente reporta una
t\'ecnica padre (p.\ ej.\ \technique{T1547}) mientras que una regla
est\'atica reporta la subt\'ecnica espec\'ifica responsable
(\technique{T1547.001}): el mismo mecanismo de persistencia
subyacente, observado en dos niveles de especificidad distintos por
dos fuentes distintas, pero que nunca se corroboran entre s\'i bajo un
agrupamiento por coincidencia exacta. Esta es una inconsistencia real,
no hipot\'etica: el propio emparejamiento a nivel de familia del
arn\'es de evaluaci\'on (\S\ref{sec:evaluation}) ya trata una t\'ecnica
padre y su subt\'ecnica como equivalentes a efectos de puntuaci\'on,
mientras que el modelo de confianza que produce el resultado puntuado
no aplicaba el mismo principio.

\code{apply\_cross\_level\_corroboration()} cierra este hueco: tras la
pasada normal de reconciliaci\'on por t\'ecnica, las t\'ecnicas se
agrupan por identificador base (quitando cualquier sufijo de
subt\'ecnica), y para cualquier grupo donde exista evidencia est\'atica
\emph{y} din\'amica entre sus miembros pero a un miembro individual le
falte una de las dos, su confianza se recalcula usando el mejor nivel
del grupo para el lado que falta, la misma regla de ``alta salvo que
ambas sean bajas'', extendida entre identificadores hermanos en lugar
de solo dentro de uno. Cada promoci\'on queda registrada en el propio
registro de la t\'ecnica (qu\'e t\'ecnica hermana aport\'o la
corroboraci\'on que faltaba), de modo que el mecanismo es auditable en
lugar de mezclar evidencia entre identificadores de forma silenciosa.
En las tres muestras validadas este cambio no alter\'o ning\'un valor
final de confianza: cada grupo con granularidad mixta actualmente
alcanzaba ya \emph{alta} de forma independiente por la propia
severidad de su fuente, o ambos miembros compart\'ian la misma \'unica
fuente. Con todo, se construy\'o \emph{antes}, no despu\'es, de la
ampliaci\'on de cobertura descrita a continuaci\'on, precisamente
porque las nuevas reglas est\'aticas espec\'ificas de subt\'ecnica
combinadas con firmas de comunidad de CAPE frecuentemente a nivel de
t\'ecnica padre hacen que este hueco sea m\'as probable en adelante,
no menos.

\subsection{El bucle de reenv\'io}
\label{sec:resubmission}

La capacidad real del malware multietapa suele distribuirse entre
varios binarios: RoningLoader, por s\'i solo, a lo largo de su cadena
de infecci\'on suelta o carga reflexivamente un driver rootkit en modo
kernel, una DLL cuya \'unica funci\'on es enumerar y terminar productos
de seguridad, y un cliente troyano de acceso remoto independiente
(\S\ref{sec:case-roning}). Fusionar ingenuamente la evidencia de cada
componente soltado en el propio \code{attck\_mapping.json} de la
muestra padre atribuir\'ia mal la confianza: el modelo de
\code{reconcile.py} asume que toda observaci\'on agrupada bajo una
t\'ecnica describe el \emph{mismo} binario, y una carga soltada no
demuestra que el cargador padre presente ese comportamiento.

El bucle de reenv\'io, en cambio, hace pasar a cada componente
soltado/extra\'ido por su propia ejecuci\'on completamente
independiente a trav\'es de la misma pipeline de an\'alisis
est\'atico $\rightarrow$ normalizaci\'on $\rightarrow$ mapeo, etiquetada
con linaje expl\'icito (hash del padre, mecanismo de descubrimiento,
marca de tiempo de descubrimiento) hacia la muestra que la produjo.
Tres componentes, divididos por dependencia y no por conveniencia:

\begin{enumerate}
  \item \textbf{Escritor de manifiesto} (se ejecuta en el host, solo con
    biblioteca est\'andar): localiza los bytes de ficheros soltados en
    el propio almacenamiento de resultados de CAPE, los copia a una
    cola, y escribe un manifiesto de linaje por candidato, acotado por
    un l\'imite configurable de artefactos.
  \item \textbf{Procesador est\'atico} (se ejecuta dentro del contenedor
    de an\'alisis est\'atico, que ya tiene las dependencias m\'as
    pesadas, \code{pefile}, \code{yara-python}, FLOSS, Ghidra, que
    necesita el an\'alisis de ficheros soltados): analiza cada artefacto
    en cola y etiqueta su informe con los datos del manifiesto de
    linaje. Duplicar esas dependencias en un contenedor separado y m\'as
    ligero se consider\'o y se descart\'o por ser superficie de
    mantenimiento innecesaria sin beneficio de aislamiento.
  \item \textbf{Procesador de fusi\'on} (se ejecuta dentro del
    contenedor de la pipeline, solo con biblioteca est\'andar):
    localiza cada informe etiquetado con linaje que a\'un no se haya
    normalizado ni mapeado, y lo hace pasar por la misma pipeline de
    reconciliaci\'on de confianza que una muestra de nivel superior,
    escribiendo en una ruta de salida espec\'ifica del componente,
    nunca en el propio \code{attck\_mapping.json} de la muestra padre,
    de modo que la evidencia de un componente soltado nunca pueda
    sobrescribir o mezclarse silenciosamente con el propio resultado
    puntuado del padre. Esto impone a nivel de sistema de ficheros el
    mismo principio de no mezcla que \code{reconcile.py} impone a nivel
    de modelo de datos.
\end{enumerate}

Para RoningLoader en concreto, el bucle de reenv\'io descubri\'o y
analiz\'o de forma independiente 42 componentes m\'as all\'a del propio
binario cargador padre: el driver rootkit, la DLL antiantivirus y el
cliente RAT, cada uno con su propia entrada dedicada en los datos de
referencia manuales, m\'as otros 39 artefactos adicionales de carga y
configuraci\'on extra\'idos sin datos de referencia propios redactados
individualmente. La Secci\'on~\ref{sec:evaluation} punt\'ua los tres
componentes nombrados frente a sus datos de referencia por componente,
ya que puntuar solo el componente padre subestima estructuralmente lo
que encuentra la pipeline completa, incluyendo el bucle de reenv\'io;
los 23 restantes son superficie de evaluaci\'on real adicional en un
sentido distinto (\S\ref{sec:limitations}, punto~\ref{lim:n3}), cada uno
confirmado como productor de un resultado sensato, sin ca\'idas y
etiquetado con linaje, en lugar de aportar nuevos puntos de datos de
precisi\'on/exhaustividad puntuados.

\subsection{Ampliaci\'on de la cobertura de reglas de detecci\'on}
\label{sec:coverage}

Al inicio de la validaci\'on descrita en \S\ref{sec:evaluation}, cada
regla est\'atica exist\'ia porque evidencia espec\'ifica en una de las
tres muestras la justificaba. Es un enfoque s\'olido para evitar el
sobreajuste, pero tambi\'en, por construcci\'on, deja fuera buena parte
de la matriz completa de MITRE ATT\&CK: cubre lo que necesitan las
tres familias validadas, no la matriz en su conjunto. Para
cuantificar esa cobertura con precisi\'on en lugar de estimarla, se descarg\'o el
paquete STIX oficial actual de MITRE Enterprise
ATT\&CK~\cite{mitre-attack-stix} y se cont\'o directamente: 474
t\'ecnicas relevantes para Windows (incluyendo subt\'ecnicas) en la
matriz actual, de las cuales aproximadamente 30 ten\'ian una regla
est\'atica dedicada antes de esta ampliaci\'on.

Esas 474 t\'ecnicas ya son un subconjunto filtrado, no la matriz
completa: de las 697 t\'ecnicas actuales de MITRE Enterprise ATT\&CK
(excluyendo las revocadas u obsoletas), solo el 68\% est\'a etiquetado
por la propia MITRE como relevante para Windows. Adem\'as, la cobertura
a\~nadida no se reparte uniformemente entre t\'acticas: Privilege
Escalation, Persistence, Discovery, Defense Evasion, Credential Access
y Execution concentran cerca del 74\% de las etiquetas t\'actica-t\'ecnica
del conjunto de reglas, mientras que Initial Access, Lateral Movement,
Reconnaissance y Resource Development, las t\'acticas que la acotaci\'on
m\'as abajo en esta misma secci\'on argumenta que encajan mal con un
an\'alisis de detonaci\'on de un solo binario, no tienen ninguna regla
dedicada.

La cobertura se ampli\'o posteriormente a lo largo de seis lotes, de
una a dos t\'acticas cada vez, a\~nadiendo aproximadamente 55
identificadores de t\'ecnica adicionales (listado completo en el
Anexo~\ref{app:rules}). Cada adici\'on sigue la misma disciplina de
procedencia: el patr\'on de detecci\'on es el mecanismo \emph{literal}
que nombra la propia descripci\'on de la t\'ecnica en MITRE (un valor
de registro espec\'ifico, una combinaci\'on de clases WMI, un marcador
de fichero, una DLL, o un indicador de l\'inea de comandos), nunca una
API o herramienta gen\'erica por s\'i sola cuando ese nombre tambi\'en
tiene usos leg\'itimos habituales. Tras cada lote, se volvi\'o a
ejecutar el an\'alisis est\'atico en las tres muestras validadas y se
comprob\'o que no surgiera ning\'un hallazgo nuevo antes de considerar
seguro conservar el lote; ninguna de las 55 adiciones se activ\'o en
ninguna de las tres muestras, confirmando que la ampliaci\'on no
introduce nuevo riesgo de falsos positivos en el conjunto sobre el que
se miden las cifras de este proyecto. Se trata expl\'icitamente de
cobertura orientada a muestras a\'un no vistas; no mueve, ni se afirma
que mueva, las cifras de evaluaci\'on reportadas en
\S\ref{sec:evaluation}.

La ampliaci\'on de cobertura se acot\'o deliberadamente, no de forma
exhaustiva, tras comprobar el n\'umero real de t\'ecnicas para cada
t\'actica restante de MITRE (\S\ref{sec:future-work}): Reconnaissance y
Resource Development quedan estructuralmente fuera de alcance sin
importar el esfuerzo, ya que ambas describen actividad del atacante
que ocurre antes de que la muestra de malware exista; Initial Access y
Lateral Movement est\'an t\'ecnicamente dentro de alcance pero encajan
mal con lo que puede observar un an\'alisis de detonaci\'on de un solo
binario: para una familia de ransomware como Akira en concreto, el
movimiento lateral real lo suele ejecutar el operador humano con
herramientas gen\'ericas (RDP, PsExec, Cobalt Strike) entre el acceso
inicial y el despliegue del propio binario de ransomware, de modo que
la capacidad normalmente no est\'a codificada en la muestra detonada
en absoluto, un mal encaje que ning\'un banco de pruebas de sandbox
mayor cambiar\'ia.

\subsection{Despliegue y reproducibilidad}
\label{sec:deployment}

El prototipo y cada entorno de an\'alisis del que depende, no solo el
c\'odigo propio de la pipeline, est\'a definido en un \'unico
\code{docker-compose.yml}, dividido en dos perfiles activables de forma
independiente. El perfil \code{core} (motor de an\'alisis est\'atico,
pipeline, ATT\&CK Navigator) no necesita configuraci\'on a nivel de
host m\'as all\'a de Docker en s\'i, y es el que produjo todos los
resultados solo-est\'aticos de esta tesis. El perfil \code{sandbox}
(CAPE, MISP y sus servicios de apoyo: MySQL, Redis, INetSim) necesita
adem\'as una configuraci\'on manual \'unica del host (libvirt/KVM y una
m\'aquina virtual Windows invitada instant\'aneable), ya que la imagen
de disco de una VM no puede aprovisionarse con \code{docker-compose}.
La Tabla~\ref{tab:deployment} enumera cada servicio y su perfil.

\begin{table}[h]
\centering
\scriptsize
\begin{tabular}{@{}llp{6.6cm}c@{}}
\toprule
\textbf{Servicio} & \textbf{Perfil} & \textbf{Funci\'on} & \textbf{Memoria} \\
\midrule
\code{static} & core & Motor est\'atico PE/ELF (\code{pefile}, YARA, FLOSS, Ghidra headless) & 3g \\
\rowline
\code{pipeline} & core & Normalizador, mapeador, exportador (STIX/MISP) & 1g \\
\rowline
\code{navigator} & core & Visualizaci\'on de capas del ATT\&CK Navigator & 512m \\
\rowline
\code{redis} & sandbox & Backend de sesi\'on/trabajos de MISP & 256m \\
\rowline
\code{inetsim} & sandbox & Red simulada para la detonaci\'on, sin salida real & 256m \\
\rowline
\code{cape} & sandbox & Controlador del sandbox CAPE (detonaci\'on KVM/libvirt) & 6g \\
\rowline
\code{mysql} & sandbox & Almacenamiento relacional de MISP & 1g \\
\rowline
\code{misp} & sandbox & N\'ucleo de MISP (almacenamiento de eventos, API, datos de galaxia) & 2g \\
\bottomrule
\end{tabular}
\caption{Inventario de servicios por perfil, a partir de \code{docker-compose.yml}.}
\label{tab:deployment}
\end{table}

El \'unico paso de configuraci\'on que \code{docker-compose}
genuinamente no puede realizar es aprovisionar el hipervisor
KVM/libvirt y su red virtual en el propio host, ya que \code{cape} e
\code{inetsim} necesitan un \code{/dev/kvm} real y un puente libvirt
real al que conectarse, no un namespace de contenedor.
\code{scripts/host-prereqs.sh} automatiza este paso \'unico y es
idempotente: instala \code{qemu-kvm} y \code{libvirt}, verifica que
\code{/dev/kvm} sea realmente utilizable, y descubre la interfaz
puente, la subred y la IP de puerta de enlace que libvirt est\'a usando
en esa m\'aquina concreta (deliberadamente no codificado de forma fija,
ya que var\'ia de host a host), escribiendo los valores descubiertos en
\code{.env} para que \code{inetsim} y la propia interpolaci\'on de
configuraci\'on de CAPE puedan recogerlos sin reconstruir la imagen.

Dos decisiones de dise\~no m\'as en el fichero de composici\'on se derivan
directamente de la postura de contenci\'on prioritaria descrita en
\S\ref{sec:dynamic}. Primero, los servicios de backend (an\'alisis
est\'atico, pipeline, MySQL, Redis) se sit\'uan en una red Docker con
\code{internal: true} y sin ninguna ruta hacia el host ni hacia
internet; solo MISP y el ATT\&CK Navigator, los dos servicios que un
humano necesita alcanzar desde un navegador, se sit\'uan en una segunda
red no interna, manteniendo el l\'imite de no salida en el punto m\'as
estrecho que todav\'ia permite a un humano revisar los resultados.
Segundo, CAPE e INetSim se ejecutan con red de host en lugar de en la
red de backend aislada, una excepci\'on deliberada y acotada: la
detonaci\'on de CAPE basada en KVM/libvirt no tiene ning\'un patr\'on de
despliegue Docker oficialmente soportado, y todo enfoque de la
comunidad converge en red de host para el propio controlador del
sandbox, de modo que el l\'imite de no salida se impone una capa m\'as
abajo, en la propia configuraci\'on de red de la VM invitada, en lugar
de en la capa Docker para este servicio concreto. Ejecutar ambos
perfiles juntos en un host de 16GB se acerca a unos 14GB, algo que solo
encaja porque la detonaci\'on es secuencial, una VM invitada cada vez,
nunca concurrente.
